%%% DRAFT NOTES (v1, 2026-10-16)
%%%
%%% Modeled on "Executable Archaeology: Reanimating the Logic Theorist"
%%% (Moews & Shrager, IEEE Annals, lt-v2-revised-6.tex): same class,
%%% same section pattern (history / reanimation A / reanimation B /
%%% what it shows / what executable archaeology contributes / related
%%% work / next steps / personal coda).
%%%
%%% Sources for every factual claim about the two reanimations:
%%%   ../repo/HeuristicCoder/HeuristicCoder_guide.md, seshsums/, run/
%%%   ../KBEmacs/repo/hackers-guide.md, hackers-guide-sqrt.md,
%%%     seshsums/, code/TRANSCRIPTION-FIXES.md, code/volume_1/358-test.lisp
%%% Quotations from Simon (1963) and Waters (1986) were checked against
%%% text extracted from the PDFs in those directories.
%%%
%%% TODO (JS):
%%%   - Venue/department: kept the Annals class and "ANECDOTES" header
%%%     from the LT paper; this may be better as a feature article.
%%%   - Authorship: should Kent Pitman (source of the KBE listings, and a
%%%     KBEmacs team member) be invited as co-author, or at least asked
%%%     to read the KBEmacs section? Same question for anyone at CMU.
%%%   - Check the characterization of modern AI coding tools; it is kept
%%%     deliberately generic (no product claims). Add citations if the
%%%     editors want specific systems named.
%%%   - Verify all bibliography entries marked "TODO verify".
%%%   - The coda is written in the first person, as in the LT paper.
%%%   - KBEmacs: if the other five 1984 sequences (BUCKET, DELETE, HASH,
%%%     MEMQ, SUMLIST) get run before submission, update the KBEmacs
%%%     section and "Next Steps". Reproducing the three 1984 *failures*
%%%     would be the strongest fidelity result in the paper.
%%%   - Heuristic Compiler: the 1961 natural-language front end (T50-T300,
%%%     U23-U99, U199) has not been run; likewise update if it is.

\documentclass{IEEEcsmag}

\usepackage[colorlinks,urlcolor=blue,linkcolor=blue,citecolor=blue]{hyperref}
\expandafter\def\expandafter\UrlBreaks\expandafter{\UrlBreaks\do\/\do\*\do\-\do\~\do\'\do\"\do\-}
\usepackage{upmath,color}

\jvol{XX}
\jnum{XX}
\paper{PP}
\jmonth{Month}
\jname{\textit{IEEE Annals of the History of Computing}}
\jtitle{\textit{IEEE Annals of the History of Computing}}
\pubyear{YYYY}

\usepackage{url}
\usepackage{booktabs}
\usepackage{listings}
\lstset{basicstyle=\footnotesize\ttfamily,columns=fullflexible,keepspaces=true}

\setcounter{secnumdepth}{0}

\begin{document}

\sptitle{DEPARTMENT: ANECDOTES}

\title{Before the Copilots: Reanimating Simon's Heuristic Compiler (1961) and MIT's KBEmacs (1984)}

\author{Jeff Shrager}
\affil{Bennu Climate, Inc., Palo Alto, CA 94301 USA}

\markboth{ANECDOTES}{ANECDOTES}


\begin{abstract}
AI-assisted programming is often presented as a new phenomenon, but
the idea that a machine might write or help write programs is as old
as artificial intelligence itself. We describe the reanimation of two
historical AI programming systems from their original source listings.
The first is Herbert Simon's Heuristic Compiler (1961), an IPL-V
program that treated writing a program as a problem to be solved by
the means-end analysis of the General Problem Solver. The second is
KBEmacs (1984), the Knowledge-Based Editor of MIT's Programmer's
Apprentice project, which let a programmer build programs by combining
standard algorithmic ``clich\'es'' through a stylized English command
language. Both run again: the Heuristic Compiler reproduces the
compiled code printed in Simon's 1963 paper, and KBEmacs reproduces one
of its own December 1984 self-tests exactly. Together they illustrate
two historical approaches---autonomous problem solving and the
knowledgeable apprentice---whose concerns reappear, in different
technical form, in today's AI coding assistants.
\end{abstract}

\maketitle

\begin{IEEEkeywords}
artificial intelligence history,
automatic programming,
AI-assisted programming,
digital preservation,
Heuristic Compiler,
IPL-V (Information Processing Language V),
KBEmacs,
Lisp Machine,
Programmer's Apprentice,
H.~A.~Simon,
software archaeology,
R.~C.~Waters
\end{IEEEkeywords}

\section{INTRODUCTION}
\label{sec:intro}

Programmers in the 2020s increasingly work alongside AI systems that
complete, write, explain, test, and repair code. These systems, built
on large language models trained on enormous corpora of source
code~\cite{chen2021codex,jimenez2024swebench}, are widely described as
something new under the sun. In one sense they are: nothing like their
scale, fluency, or breadth existed before. But the \emph{aspiration}
they embody, a machine that takes part in the work of programming, is
nearly as old as programming. In the 1950s ``automatic programming''
meant what we now call compiling: FORTRAN was introduced as ``the
FORTRAN automatic coding system''~\cite{backus1957fortran}. Within a few
years of the founding of artificial intelligence, researchers were
asking whether the methods of AI could be turned on the programming
task itself~\cite{balzer1985,richwaters1988myths}.

This paper describes the reanimation, from their original source
listings, of two historical systems that represent two quite different
answers to that question.

The first is Herbert A.~Simon's \emph{Heuristic Compiler}, written in
IPL-V in 1961 and described in a RAND report that year~\cite{simonP2349}
and in the \textit{Journal of the ACM} in 1963~\cite{simon1963jacm}.
Simon asked ``whether the task of writing a computer program can be
regarded as a `problem' in the sense in which that term is used in
GPS,'' the General Problem Solver of Newell, Shaw, and
Simon~\cite{newell1959gps}. The Heuristic Compiler is, in effect, GPS
turned on programming: given a description of what a routine should do,
it searches for a program by finding differences between what it has
and what it wants, and applying operators relevant to those
differences.

The second is \emph{KBEmacs}, the Knowledge-Based Editor in Emacs,
built by Richard C.~Waters with Kent Pitman, Charles Rich, and others at
the MIT Artificial Intelligence Laboratory as the culminating
demonstration system of the Programmer's Apprentice
project~\cite{waters1985tr,waters1986aimag,richwaters1987overview}.
KBEmacs did not try to write programs on its own. It was an
\emph{assistant}: the programmer chose which standard algorithms
(``clich\'es'') to combine, and told KBEmacs so in a stylized English
command language; KBEmacs supplied the details and wrote, and
continually rewrote, the code.

The two systems are separated by a quarter century, by programming
language (IPL-V on 1960s mainframes versus Zetalisp on the Symbolics
3600 Lisp Machine), and by research tradition (Carnegie Tech and RAND
versus MIT). Yet between them they stake out most of the conceptual
territory that today's AI coding tools occupy: how a programming task
is specified, where the knowledge needed to carry it out comes from,
how work is divided between human and machine, and how the system
checks and explains what it has done.

We approached both as exercises in what has been called
\emph{executable archaeology}~\cite{shrager2026,moewsshrager2026}:
studying historical computational systems by reconstructing and
running their original code, rather than reimplementing their
algorithms. The Heuristic Compiler was transcribed card by card from a
1961 line-printer listing in Simon's papers and run on a new IPL-V
interpreter; KBEmacs was transcribed from a four-volume January 1985
printout and run under a new Zetalisp compatibility layer for a modern
Common Lisp. In both cases the reanimation surfaced details invisible
in the published accounts. And, in a turn that would surely have
amused both Simon and Waters, both reanimations were themselves carried
out with substantial help from a modern AI coding assistant.

\section{HISTORICAL BACKGROUND}
\label{sec:history}

\subsection{From Automatic Coding to Heuristic Programming}

The phrase ``automatic programming'' has meant different things in
different decades, usually whatever was just beyond the current state
of the art~\cite{balzer1985,richwaters1988myths}. For the designers of
FORTRAN it meant translating algebraic formulas into machine
code~\cite{backus1957fortran}. Such translators are deterministic: the
programmer has already decided \emph{how} the computation is to be done,
and the compiler maps that decision into another notation.

Newell, Shaw, and Simon's work at RAND and Carnegie Tech in the 1950s
took a different view of what computers could do. Their Logic Theorist
(1956) searched for proofs of theorems in \textit{Principia
Mathematica}~\cite{LT1957,moewsshrager2026}, and their General Problem
Solver (1959) generalized its methods into \emph{means-end
analysis}~\cite{newell1959gps}. Simon's 1963 summary of the method is
worth quoting in full, because the Heuristic Compiler implements it
almost literally:

\begin{quote}
Given the present and desired objects, find a difference between them.
Next, find an operator relevant to the difference; determine if the
operator can be applied to the present object. If so, apply it. (If
not, describe the objects to which it would apply and transform the
present object into an object of that kind---a new ``desired
object.'') Take the new object thus obtained as the present object and
repeat the process.~\cite[p.~494]{simon1963jacm}
\end{quote}

Simon then made the move that defines the Heuristic Compiler. A proof
is a sequence of operators that transforms axioms into a theorem; read
the other way, it is a program that generates the theorem from the
axioms. ``Thus a theorem-proving system can be regarded, at least
formally, as a program-writing system. Conversely, if we can formulate
a programming goal as a difference between a present and a desired
object, we can presumably use the same processes \ldots\ to generate a
program''~\cite[p.~495]{simon1963jacm}. The same observation, in a
deductive rather than heuristic setting, later became the foundation of
program synthesis from theorem
proving~\cite{green1969,mannawaldinger1971}.

Simon described the program in RAND report P-2349 (June 30, 1961),
which included ``a listing of the program as of June 1961''
~\cite{simonP2349}; in the \textit{Journal of the ACM} in
1963~\cite{simon1963jacm}; and at greater length in RAND Memorandum
RM-3588-PR, \textit{The Heuristic Compiler} (May 1963), whose
appendices list a 1963 version of the compiler and a small companion
program for ``annexing'' information stated as definite
descriptions~\cite{simonRM3588}. He was explicit about its status:
``no claims will be made for the heuristic programming procedures
described here as practical approaches to the construction of
compilers. Their interest lies in what they teach us about the nature
of the programming task''~\cite[p.~493]{simon1963jacm}.

\subsection{The Programmer's Apprentice}

The Programmer's Apprentice (PA) project began at MIT in the
mid-1970s, with Charles Rich and Howard Shrobe's \textit{Initial Report
on a LISP Programmer's Apprentice}~\cite{richshrobe1976}. Its premise
was that fully automatic programming of realistic software was a
distant goal, and that a more productive near-term aim was an
intelligent \emph{assistant} that would share with an expert
programmer a body of knowledge about programs. Waters later put it
this way: ``Recognizing that it will be a long time before it is
possible to fully duplicate an expert programmer's abilities, the
project seeks to develop an intelligent assistant system \ldots\ which
will help a programmer in various phases of the programming
task''~\cite[p.~47]{waters1986aimag}.

Over the following decade the project developed three central ideas.
\emph{Clich\'es} are standard computational methods, such as a linear
search, a report generator, or successive approximation, each with
\emph{roles} that vary from use to use and a fixed surrounding
structure. \emph{Plans} are a language-independent representation of
programs as hierarchical flowcharts with explicit data and control flow,
which Rich formalized as the Plan
Calculus~\cite{richwaters1987overview}. And the \emph{assistant
approach} itself: a division of labor in which the programmer makes the
design decisions and the apprentice handles the routine details.
Waters's dissertation developed the analysis of programs into
plans~\cite{waters1978}.

KBEmacs was the realization of these ideas as a working editor. It ran
on Symbolics Lisp Machines inside the Zmacs editor, and handled both
Lisp and Ada. Its best-known demonstration builds an Ada report
generator, \texttt{UNIT-REPAIR-REPORT}, from five commands such as
``Define a simple-report procedure UNIT-REPAIR-REPORT'' and ``Fill the
enumerator with a chain-enumeration of UNITS and
REPAIRS''~\cite[p.~48]{waters1986aimag},\cite{waters1985tse}. Waters
pointed to the order-of-magnitude difference between the size of the
command sequence and the size of the resulting program as a measure of
the leverage KBEmacs provided, while noting that program construction
is only a small fraction of the software life
cycle~\cite[p.~48]{waters1986aimag}.

KBEmacs was candidly a research prototype. By 1986 Waters reported that
``knowledge-based operations on large programs can take longer than 5
minutes,'' that the system handled ``only about 50\% of Ada'' and 90\%
of Lisp, that it knew ``only a few dozen clich\'es'' where a practical
tool ``would need to know many hundreds,'' and that, its ideas having
been used ``for essentially all [they are] worth,'' KBEmacs ``has been
mothballed''~\cite[pp.~54--55]{waters1986aimag}. The project continued
with work on the Plan Calculus, on automated reasoning, and on
requirements and design apprentices, culminating in Rich and Waters's
1990 book~\cite{richwaters1990book}.

\begin{table*}[t]
\centering
\caption{The two reanimated systems.}
\label{tab:systems}
\footnotesize
\begin{tabular}{@{}p{0.18\textwidth}p{0.36\textwidth}p{0.4\textwidth}@{}}
\toprule
 & \textbf{Heuristic Compiler} & \textbf{KBEmacs} \\ \midrule
Principal author(s) & H.~A.~Simon (with H.~S.~Kelly) & R.~C.~Waters, K.~Pitman, C.~Rich, and others \\
Institution & Carnegie Tech / RAND & MIT AI Laboratory \\
Surviving source & Listing ``HEURISTIC CODER 7/16/61'' & Four-volume printout, January 1985 (KBE 358.0x) \\
Language / platform & IPL-V / (machine not established) & Zetalisp / Symbolics 3600, Release 4.5 \\
Size & about 4,470 cards & 87 files, about 34,500 lines transcribed \\
Input & state or functional description & stylized English commands \\
Central method & means-end analysis (GPS) & instantiating and combining clich\'es in plans \\
Reanimated on & new Common Lisp IPL-V interpreter & Zetalisp compatibility layer for SBCL \\
Benchmark reproduced & code printed in Simon (1963) & self-test of December 14, 1984 \\ \bottomrule
\end{tabular}
\end{table*}

\section{REANIMATING THE HEURISTIC COMPILER}
\label{sec:hc}

\subsection{The Listing}

The surviving source is a line-printer listing headed
\texttt{HEURISTIC CODER 7/16/61}, preserved in the Herbert A.~Simon
Papers at the Carnegie Mellon University
Archives~\cite{simonlisting1961}. It was located by Mia Golec and Emily
Davis of CMU with the author. Dated two weeks after RAND report
P-2349, it is very nearly the program that report describes. The 95
scanned pages print the complete card deck, one card per line, with
comments, IPL fields, and card identifiers. One page was scanned twice;
the remaining 94 hold 4,473 cards, close to the pencilled total of
``4464'' on the first page. The listing also carries pencil:
check marks, stack-depth counts beside a few routines, and a handful of
proposed corrections, apparently the original debuggers' notes.

The deck turned out to contain three things: an executive and two
compilers; a natural-language front end (a character scanner,
dictionary, and parser) intended to let the compiler accept English
definitions; and a small ``world'' of knowledge about IPL-V, encoded as
data, that the compilers reason about. The listing ends with the data
and has no start card, so the program was evidently started by a
separate card or deck.

\subsection{Everything Is a Description List}

The central representational idea of the Heuristic Compiler is that a
\emph{routine} is an object with attributes, stored as an IPL-V
description list (a property list, in later Lisp terms). Each attribute
is one possible definition of the routine: its IPL name, its
\emph{functional description} (a process name plus arguments, such as
``INSERT ( ) AT END OF ( )''), its \emph{state description} (which cells
it affects, and what each holds before and after), its flow diagram,
and its IPL-V code. Below that level, an IPL instruction is itself a
description list with attributes for each field, and a symbol such as
\texttt{J3} is a description list with a region (the letter J) and a
location (the integer 3). The program never handles IPL code as text;
it builds code out of these structures and prints it only at the end.
Simon remarked that description lists gave IPL-V ``expressive
capabilities not readily available in more usual computer command
languages''~\cite[p.~506]{simon1963jacm}.

The executive (routines U135--U138) is the paper's general compiler,
almost card for card. It finds the ``closest'' definition the routine
has, looks up in a table of connections the compiler relevant to that
kind of definition, applies it, and tests whether the routine has moved
to a better kind of definition. If not, it gives up. It is a small
means-end loop whose \emph{differences} are kinds of definition and
whose \emph{operators} are whole compilers.

The \emph{state description compiler} (U140--U169 and X5--X7) is the
most GPS-like part. It compares the input and output states of each
affected cell, classifies the biggest difference (addition, deletion,
replacement, or copy), looks up an operator relevant to it, binds the
operator's variables to the actual cells, and then recursively compiles
two subproblems: getting from the original input state to the state the
operator requires, and getting from the operator's output state to the
desired output. Its knowledge of IPL-V consists of exactly two
operators, LOAD and REPLACE, each itself a tiny routine with its own
state description.

The \emph{functional description compiler} (U134, U130, and helpers)
works by analogy. Given a routine described as, say, ``INSERT (1) AT
END OF THE VALUE OF THE ATTRIBUTE (0) OF (2),'' it searches the routines
already compiled for one with the same \emph{process}, here ``INSERT (0)
AT END OF (1),'' whose code is the single primitive \texttt{J65}. It
then adapts that routine to the new arguments: it moves inputs into
working cells, compiles code to fetch each argument, recursing when an
argument is itself a described computation (a ``determiner''), and
wraps the result in instructions that save and restore the working
cells.

\subsection{Running It}

The deck runs on a Common Lisp IPL-V interpreter originally written to
run the Logic Theorist~\cite{shrager2026,shragerrepo}, extended to read
80-column card images directly and to supply J-functions the Heuristic
Compiler uses but the Logic Theorist did not. Several existing
J-functions turned out to be wrong in ways the Logic Theorist never
exercised; after they were corrected against the 1964
manual~\cite{newell1964ipl}, the interpreter was rechecked and still
reproduces Stefferud's 1963 Logic Theorist output for all 24 of its
theorems~\cite{earepo}.

One print routine, U126, is called but absent from the 1961 listing; it
was supplied from the 1963 listing in RM-3588-PR,
Appendix~A~\cite{simonRM3588}, and affects only printing. With that and
a start card for the demonstration routine T1 (the 1963 listing ends
with exactly such a card), the program runs Simon's published examples.

\begin{figure*}
\begin{lstlisting}
 FIND THE VALUE-OF THE ATTRIBUTE (0) OF (1)              <- an already-compiled routine

                                       J10   0

                                       J65   0           <- the code of INSERT (0) AT END OF (1)

                                     11W1                <- fetch argument (1), now in W1
                                       J65   0

                                     11W2                <- fetch the determiner's arguments
                                     11W0                   and apply its code (J10)
                                       J10
                                     11W1
                                       J65   0

 J13
 INSERT (1) AT END OF  THE VALUE-OF THE ATTRIBUTE (0) OF (2)

                               J13     J52               <- preserve working cells W0-W2
                                     11W2
                                     11W0
                                       J10
                                     11W1
                                       J65
                                       J32   0           <- restore them and exit
\end{lstlisting}
\caption{The 1961 Heuristic Compiler, running in 2026, compiling INSERT
AT END OF VALUE LIST from its functional description. The intermediate
stages are printed by the program itself (a print in routine U113 that
appears to be a debugging aid left in the deck). The final code is
character for character the code printed in Simon's 1963 paper.
Annotations after \texttt{<-} are ours.}
\label{fig:insert}
\end{figure*}

\textbf{Compiling from a functional description.} Given the functional
description of INSERT AT END OF VALUE LIST, the program produces
\texttt{J13 J52 11W2 11W0 J10 11W1 J65 J32 0}, character for character
the code in the 1963 paper~\cite[pp.~499--500]{simon1963jacm}
(Figure~\ref{fig:insert}). A print left in routine U113 shows the code
growing one argument at a time, which happens to reproduce the sequence
of stages through which the paper walks its readers. We note one
qualification. In Simon's runs, as in ours, the functional description
is supplied as a hand-coded list structure, not as typed English; the
phrase is printed from it, but the compiler matches processes by symbol
identity. The English front end in the deck has not yet been run (see
``Next Steps'').

\textbf{Compiling from a state description.} Simon's running example of
the state description compiler is J3, SET SIGNAL MINUS: a routine that
replaces the top symbol of the signal cell H5 with the symbol J3
(``minus''). The state description says that H5 goes from \emph{some
symbol on top of the rest of its pushdown list} to \emph{J3 on top of
the same rest}. The compiler finds a replacement difference in a named
cell; finds REPLACE relevant, with its variable bound to H5, giving
\texttt{20H5}; sets up the prefix subproblem ``get J3 into the
communication cell,'' an addition, for which LOAD is relevant, giving
\texttt{10J3}; finds no suffix subproblem; and assembles
\texttt{10J3 20H5 0}, the code in the paper~\cite[p.~498]{simon1963jacm}.

There is a catch, and it is the kind of thing only running the program
reveals. \emph{As listed, the deck does not compile J3.} The listing's
J3 has both a functional and a state description, and the executive
always prefers the functional description, on the stated assumption
that ``it is easier to compile from a functional description than from
a state description''~\cite[p.~502]{simon1963jacm}. No already-compiled routine shares J3's process,
the functional description compiler fails, the executive sees no
progress, and it gives up without ever trying the state description.
The paper's J3 result is obtained only when J3's functional description
is removed before the run, after which the state description compiler
produces exactly the published code. The 1963 listing in RM-3588-PR
suggests what happened: its executive is card for card the same, but
the cards that give J3 a functional description are numbered as later
insertions. Most likely the paper's J3 came from a version in which J3
had only a state description. (The 1963 state description compiler is
also simpler in its representation of states, holding bare symbols
where the 1961 version holds described ones.)

\textbf{Flow diagrams.} The paper's last section handles routines with
branches and loops by decomposing them into \emph{segments} joined by a
flow diagram. The program reassembles routine J77 (``test whether there
is a symbol equal to 0ACCUMULATOR on list 1ACCUMULATOR'') from its
segments, and decomposes it again into the paper's flow diagram.

\textbf{The information annexer.} RM-3588-PR's Appendix~B, a small
program that stores statements such as ``X114 is the X33 of the X25 of
X105'' into description lists, either identifying an object already in
memory or creating one, was also transcribed (409 cards) and runs. Its
output matches the fragment of 1963 output that survives in the
memorandum~\cite{simonRM3588,earepo}.

\section{REANIMATING KBEMACS}
\label{sec:kbe}

\subsection{The Listings}

The surviving KBEmacs source is a January 1985 printout, four bound
volumes of Zetalisp listings, preserved and scanned by Kent M.~Pitman, a
member of the KBEmacs team~\cite{pitmanscans}. It covers the system as
of KBE version 358.0x (December 1984), although individual files carry
their own dates and versions. Volume~1 opens with something unusual and
extremely useful to an archaeologist: the results of KBEmacs's
automatic self-test, run ``Friday the fourteenth of December, 1984;
2:25:23 pm'' on the Lisp Machine \emph{Merlin}. The self-test runs six
command sequences, compares the code KBEmacs generates with expected
code, and reports that three sequences (HASH, SQRT, and SUMLIST)
succeeded and three (BUCKET, DELETE, and MEMQ) failed with ``Bad code
generated.''

The 87 files were transcribed with the assistance of a large language
model into about 34,500 lines of source~\cite{kbemacsrepo}, roughly
consistent with Waters's description of the system as ``some 40,000
lines of Lisp code''~\cite[p.~55]{waters1986aimag}. The transcription
was initially unchecked; it was checked chiefly by \emph{executing}
it.

\subsection{Text and Plans}

A KBEmacs user works in a stylized English command language phrased in
terms of clich\'es and their roles. The clich\'e library is written as
code with holes, as in the definition of successive approximation
in Figure~\ref{fig:cliche}.

\begin{figure*}
\begin{lstlisting}
(DEFINE-CLICHE SUCCESSIVE-APPROXIMATION
   (PRIMARY-ROLES (INITIAL-APPROXIMATION)
    DESCRIBED-ROLES (INITIAL-APPROXIMATION CRITERION IMPROVEMENT-STEP)
    COMMENT "iteratively improves an approximation until {the criterion} is
             satisfied" ...)
  (LET* ((APPROXIMATION {the input initial-approximation}) )
    (LOOP DO
      (IF ({the criterion} APPROXIMATION) (RETURN APPROXIMATION) )
      (SETQ APPROXIMATION ({the improvement-step} APPROXIMATION) ))))
\end{lstlisting}
\caption{A clich\'e from KBEmacs's library (\texttt{kbelib.lisp}, 1984).
The braces are KBEmacs read syntax marking \emph{roles}, the parts that
vary from use to use; the surrounding code is fixed.}
\label{fig:cliche}
\end{figure*}

But clich\'es are not stored as text. When the library loads, each
definition is run through KBEmacs's \emph{analyzer}, which turns code
into a plan: a set of facts in a small database describing segments,
ports, data flow, control flow, and roles. In a plan, a
\texttt{LOOP} with a \texttt{SETQ} disappears; what remains is a
temporal composition of a sequence of states. The knowledge-based editor
operates on plans, and after every command the \emph{coder} turns the
plan back into program text, deciding where to nest expressions and
where to introduce variables, and preferring names the programmer has
already used. Waters called this two-way translation the
\emph{representation shift}, and argued that it gave plans locality,
additivity, and modularity that text lacks~\cite{waters1986aimag}.
Because the user may also simply edit the text, KBEmacs must keep text
and plan consistent in both directions.

\subsection{Running It}

KBEmacs depends heavily on its environment: Flavors (the Lisp Machine's
object system), the Zmacs editor, Lisp Machine processes, Waters's LetS
loop notation, and Zetalisp itself. The reanimation targeted the
\emph{headless} core (the database, analyzer, clich\'e library,
command language, knowledge-based editor, coder, and explanation
generator), leaving out the editor, windows, actors, and the Ada front
end. It runs on SBCL, a modern Common Lisp, through a compatibility
layer comprising an existing Common Lisp port of Spice Lisp Flavors
extended with the Lisp Machine's \texttt{:case} method combination, a
port of Waters's LetS~\cite{g000001}, a Zetalisp reader, a library of
Zetalisp functions and conditions, and stand-ins for editor-level
facilities the core needs. A small driver replays a self-test sequence
from the 1984 record and compares the generated code with the expected
code, as the original test harness did.

The SQRT sequence (Figure~\ref{fig:sqrt}) builds a square-root routine
from three clich\'es in seven commands, and then builds a variant with
an extra parameter in seven more. After the first command that names a
clich\'e, KBEmacs holds the clich\'e's structure with its roles still
open as holes; each further command fills a role with code or with
another clich\'e, whose own unfilled roles appear as new holes. After
seven commands the plan has 32 segments and the program is complete.

\begin{figure*}
\begin{lstlisting}
>Reinitialize database.
>Define a program MYSQRT with a parameter NUM.
>Implement it as a successive-approximation.
>Implement the initial-approximation as [1.] .
>Implement the improvement-step as an average of [APPROXIMATION]
   and [(// NUM APPROXIMATION)] .
>Implement the criterion as an equality-within-epsilon of
   [(* APPROXIMATION APPROXIMATION)] and [NUM] .
>Implement the epsilon as [0.0001] .

(DEFUN MYSQRT (NUM)
  (LET* ((APPROXIMATION 1))
    (LOOP DO
      (IF (< (ABS (- (* APPROXIMATION APPROXIMATION) NUM)) 0.0001)
          (RETURN APPROXIMATION))
      (SETQ APPROXIMATION (// (+ APPROXIMATION (// NUM APPROXIMATION)) 2.0)))))

The program MYSQRT is a successive-approximation. The initial-approximation is 1.
The criterion is an equality-within-epsilon ...
\end{lstlisting}
\caption{KBEmacs, running in 2026, replaying the first half of the SQRT
sequence from its December 14, 1984 self-test. The generated code is
\texttt{equal} to the code recorded in 1984. (\texttt{//} is Zetalisp's
division operator.) The final lines are KBEmacs's own description of
the program, generated from its plan, as recorded in the 1984
self-test (abridged).}
\label{fig:sqrt}
\end{figure*}

All fourteen commands run, and both generated programs are
\texttt{equal} to the code recorded in 1984. The generated code also
works: \texttt{(mysqrt 2.0)} returns 1.4142157.

Getting there was instructive. Three kinds of problem had to be
solved, and their distribution says something about where the
difficulty of reanimation lies.

\textbf{Transcription errors.} About two dozen typo-level corrections
were needed, each checked against the scans and
logged~\cite{kbemacsrepo}. In the listing font, quote and backquote are
nearly identical, and the commonest error was one read as the other.
The most insidious were \emph{compensating} parenthesis errors: a
missing parenthesis in one place and an extra one a few lines later, so
that the file still balanced. In one case in the coder, two such pairs
silently nested about thirty function definitions inside earlier ones,
so that they were never defined. In another, in the analyzer, a clause
of a conditional migrated into the wrong branch, so that moving a
segment into a loop silently did nothing and the data-flow simplifier
retried the move forever. Neither is detectable by checking balance;
both were found by running the code and then confirmed against the
scans. In one case the transcription faithfully preserved an error in
the copy-and-paste sense: a line in the verb-definition macro had been
copied from the macro just above it, shifting every verb's arguments by
one.

\textbf{Semantics that survive only in the code.} Several behaviors of
Zetalisp in Release~4.5 had to be inferred from how KBEmacs depends on
them, because the system will not run otherwise. On the Lisp Machine,
\texttt{(listp nil)} is false. Characters behave as integers. The
function \texttt{map} maps over successive tails of a list. Most
strikingly, \texttt{(boole 4 x y)} does not mean what it means in
Common Lisp: the code is a PDP-10 Boolean-instruction code, and 4 means
``x and not y,'' which KBEmacs uses for set difference in its
live-variable analysis. With the Common Lisp meaning, the data-flow
fixed point never converges. KBEmacs's sets are kept sorted by machine
address, relying, as a comment admits, on the garbage collector never
running (the 1984 self-test header duly records ``Garbage collector
enabled: No''); addresses are emulated by creation order.

\textbf{Bugs in the reanimation's own tools.} One persistent failure
turned out to be in the modern port of LetS, which dropped the value of
its \texttt{done} form; KBEmacs uses that value to choose fresh argument
names, and without it plans acquired ports with no name.

\section{THEN AND NOW: HISTORICAL APPROACHES TO AI-ASSISTED CODING}
\label{sec:thenandnow}

Neither the Heuristic Compiler nor KBEmacs is an ancestor of today's AI
coding assistants in any direct technical sense. Large language models
are trained, not programmed with explicit operators or clich\'es, and
they operate on program text rather than on description lists or plans.
Nonetheless the questions these two systems were built to explore are
recognizably the questions that shape AI-assisted programming today.
Table~\ref{tab:compare} summarizes the comparison; we discuss five
themes.

\begin{table*}[t]
\centering
\caption{Three approaches to machine participation in programming.}
\label{tab:compare}
\footnotesize
\begin{tabular}{@{}p{0.16\textwidth}p{0.25\textwidth}p{0.25\textwidth}p{0.25\textwidth}@{}}
\toprule
 & \textbf{Heuristic Compiler (1961)} & \textbf{KBEmacs (1984)} & \textbf{LLM-based assistants (2020s)} \\ \midrule
Task specification & state or functional description & English commands naming clich\'es & natural language, code context, tests \\
Source of knowledge & hand-coded operators and routines & hand-built clich\'e library (dozens) & learned from large code corpora \\
Internal representation & description lists & plans (data and control flow) & program text (tokens) \\
Core method & means-end analysis & clich\'e instantiation, constraints & learned generation; agentic loops \\
Reuse & adapt a compiled routine with same process & instantiate a clich\'e & implicit, plus retrieval of context \\
Division of labor & autonomous & programmer decides, apprentice implements & ranges from completion to autonomy \\
Checking & progress test on kind of definition & self-test; consistency of text and plan & executing tests, human review \\
Explanation & (printed descriptions) & generated from plan & generated from text \\ \bottomrule
\end{tabular}
\end{table*}

\textbf{How is a programming task specified?} Simon's two compilers
differ precisely in how the task is described. A state description
says what the machine looks like before and after, which is close to
what we would now call a specification by examples or by
pre- and postconditions, and close to the role unit tests play in
directing a modern assistant. A functional description is an imperative
phrase, which is close to a natural-language prompt. Simon drew a
conclusion that remains apt: programmers use both, but ``at the higher
levels \ldots\ a language of functions appears generally the more
natural and powerful,'' and ``the mappings from functional descriptions
to process language are simpler and more direct than the mappings from
state descriptions''~\cite[p.~506]{simon1963jacm}. KBEmacs chose the
functional side and gave it a vocabulary of named algorithms. Today's
assistants accept both kinds of description, and much of the craft of
using them lies in combining a functional request with a test-like
check.

\textbf{Where does the knowledge come from?} Both historical systems
are knowledge-starved by modern standards. The Heuristic Compiler's
state description compiler knows exactly two operators; its functional
description compiler can compile a routine only if a routine with the
same process has already been compiled. KBEmacs knew ``a few dozen''
clich\'es; Waters estimated that a practical tool would need ``many
hundreds''~\cite[p.~55]{waters1986aimag}. Both systems identified reuse
of known solutions as the heart of the matter: Simon's functional
compiler works by finding an analogous routine and adapting it, and the
whole Programmer's Apprentice rests on the claim that most programs are
built from standard pieces. What neither could solve was acquiring that
knowledge at scale. Large language models can be seen as one answer to
precisely this bottleneck: an implicit clich\'e library absorbed from
an enormous body of code, at the price of having no explicit, inspectable
representation of what a clich\'e is or what its roles are.

\textbf{Who does what?} The Heuristic Compiler was meant to be
autonomous; KBEmacs was deliberately an apprentice, leaving design
decisions to the programmer and taking over ``the details.'' Waters's
argument for the apprentice was practical: full automation was not
near, and a semiexpert assistant could be useful at once. Modern tools
have traversed the same axis in a few years, from completing a line, to
writing a function on request, to agents that plan, edit, and test
across a code base with varying degrees of
supervision~\cite{jimenez2024swebench}. The question of where the
programmer's decisions end and the machine's begin, central to the
Programmer's Apprentice, is again a live design question.

\textbf{How does the system know it is making progress?} The Heuristic
Compiler's executive is a loop: apply a compiler, test whether the
routine has moved to a better kind of definition, and stop if it has
not. Its state description compiler is difference reduction in the
strict GPS sense. Modern coding agents commonly run a loop of a similar
shape: make a change, run the tests or the compiler, examine the
difference between what happened and what was wanted, and try again.
The difference detector is now an executed test suite rather than a
hand-coded comparison of states, and the operator is a learned model
rather than a table of connections, but the architecture of
means-end iteration is recognizable. KBEmacs, for its part, came with a
self-test in which the expected output was code, which is exactly what
made it possible to check the reanimation in 2026.

\textbf{Representation and explanation.} The two historical systems
never treated code merely as text. The Heuristic Compiler builds IPL
instructions out of description lists and prints them only at the end;
KBEmacs reasons about plans and regenerates text after every command,
taking pains to keep text and plan consistent when the user edits the
text directly. That explicit representation is what allowed KBEmacs to
describe a program in terms of its design (``The program MYSQRT is a
successive-approximation \ldots''). Today's assistants operate on text
and also explain code fluently, but their explanations are generated
from text, not read off an explicit representation of the design. The
trade-off between the robustness of working with text, which KBEmacs
found painfully expensive to analyze---Waters reports that ``well over
half of all the effort expended on KBEmacs went into implementing an
analyzer and coder''~\cite[p.~55]{waters1986aimag}---and the
inspectability of working with plans, has not gone away.

\section{WHAT EXECUTABLE ARCHAEOLOGY CONTRIBUTES}
\label{sec:contribution}

As with the Logic Theorist~\cite{moewsshrager2026}, it is fair to ask
what running these programs teaches that reading them and their
documentation does not. Four answers stand out.

\textbf{It reveals gaps between the published demonstration and the
surviving program.} The Heuristic Compiler's J3 example does not
compile from the deck as listed; the published result requires a
version of the data in which J3 lacks a functional description. This
is not visible from the paper, nor from reading the executive, whose
preference ordering looks perfectly sensible on the page. Similarly,
the ``functional descriptions'' in Simon's demonstrations turn out to
be hand-coded list structures; the English front end, which occupies a
large part of the deck, is not exercised by any of the published
examples.

\textbf{It recovers semantics that survive only in the artifacts.} The
PDP-10 meaning of \texttt{boole}, the falsity of \texttt{(listp nil)},
and the integer nature of characters on the Lisp Machine are all
documented somewhere, but nothing about KBEmacs's listings tells a
reader that its correctness depends on them; the system simply does not
work otherwise. On the IPL-V side, several J-functions had to be
corrected in ways the Logic Theorist had never exercised, so each new
program tests the interpreter as much as the interpreter tests the
program.

\textbf{It checks the transcription, and the record.} Compensating
parenthesis errors are invisible to every static check short of
re-reading the scan; executing the code finds them. Conversely, the
fact that a fresh transcription produces the Heuristic Compiler's
published code character for character, or KBEmacs's 1984 output
\texttt{equal} for \texttt{equal}, is strong evidence that the
transcription and the interpreter are right along those paths, since
errors would almost certainly have produced different code.

\textbf{It makes historical claims checkable.} Simon's 1963 claims
about what the Heuristic Compiler did, and the 1984 self-test's claims
about what KBEmacs did, can now be tested and extended. Most
interestingly, KBEmacs's self-test records three \emph{failures}. A
reanimation that reproduces those failures, with the same wrong code,
would be a sharper test of fidelity than one that reproduces only the
successes.

\section{RELATED WORK}
\label{sec:related}

This work belongs to a series of reanimations of early AI programs from
their original sources. Lane et al.\ restored Weizenbaum's ELIZA to
operation on a reconstructed CTSS running on an emulated IBM
7094~\cite{lane2025eliza}. Shrager~\cite{shrager2026} and Moews and
Shrager~\cite{moewsshrager2026} reanimated two versions of the Logic
Theorist, one in IPL-I and one in IPL-V; the IPL-V interpreter used for
the Heuristic Compiler descends from that work.

The history of automatic programming has been surveyed by
Balzer~\cite{balzer1985} and by Rich and
Waters~\cite{richwaters1988myths}, the latter from the perspective of
the Programmer's Apprentice. Deductive program synthesis, which
formalizes the identification of proofs with programs that Simon made
informally, begins with Green~\cite{green1969} and Manna and
Waldinger~\cite{mannawaldinger1971}. Within the Programmer's Apprentice
project, KBEmacs's capabilities and limits are assessed by Waters
himself~\cite{waters1986aimag} and by Rich and
Waters~\cite{richwaters1987overview,richwaters1990book}; Pitman
describes its interface~\cite{pitman1983}.

Reimplementing either system's ideas in a modern language would be a
quite different undertaking from the present work. Both reanimations
run the original code, transcribed from the surviving listings, on
reconstructions of the abstract machine (IPL-V) or language
environment (Zetalisp) it expected. For KBEmacs, a fuller reanimation
would run the complete system, including its editor, on an emulated
Lisp Machine; we have not attempted this.
%%% TODO (JS): check whether a Symbolics Release 4.5 world exists that
%%% could run under an emulator, and whether to say more here.

\section{NEXT STEPS}
\label{sec:next}

For the Heuristic Compiler, the largest piece of unexplored territory
is the English front end: the scanner, dictionary, and parser
(routines T50--T300 and U23--U99) and the routine that converts a parse
into a functional description. Simon credited H.~S.~Kelly with
``linguistic components for the compiler which are omitted from this
report''~\cite[p.~506]{simon1963jacm}; these routines may be them. Running
them requires input sentences, which are not in the listing. The pencil
corrections in the listing also remain to be evaluated.

For KBEmacs, the next step is to run the other five sequences of the
1984 self-test, aiming to reproduce both the three successes and the
three failures; then the Ada front end, which would allow the
\texttt{UNIT-REPAIR-REPORT} demonstration of Waters's papers to be
rerun.

The transcriptions, interpreters, compatibility layer, and Hacker's
Guides to both programs are open source~\cite{earepo,kbemacsrepo}.

\section{A PERSONAL CODA}

{\em Jeff Shrager}

%%% TODO (JS): revise to taste; facts are from the seshsums.
Both reanimations described here were carried out in collaboration with
a modern AI coding assistant. It drafted the first transcriptions from
the page images, wrote much of the KBEmacs compatibility layer and the
IPL-V interpreter extensions, chased down the PDP-10 meaning of
\texttt{boole}, and found the compensating parentheses. It also made
mistakes, some of which took a while to find, and none of the results
reported here were accepted until the program had reproduced its own
historical output. In short, the division of labor was very much the
one Waters proposed: I made the decisions about what to do and what
counted as success, and the apprentice did much of the detailed work.

There is a pleasing circularity in this. Simon wondered whether
programming could be treated as problem solving; Waters wondered
whether a programmer could be given an apprentice that shared knowledge
of standard algorithms. Forty and sixty-five years later, an apprentice
of a kind neither of them imagined helped bring their programs back to
life. Simon closed his paper by saying that the Heuristic Compiler was
``an exploratory tool that is beginning to reveal to us something of the
role played by these and other representations in the programming
process''~\cite[p.~506]{simon1963jacm}. Running it again, with
help, suggests the exploration is not over.

\section{ACKNOWLEDGMENTS}

%%% TODO (JS): add others as appropriate.
KBEmacs could not have been reanimated without Kent M.~Pitman, who
preserved the 1985 listings and made the scans
available~\cite{pitmanscans}. The Heuristic Coder listing was located in
the Herbert A.~Simon Papers at the Carnegie Mellon University Archives
by Mia Golec and Emily Davis. David Moews's IPL-V Logic Theorist deck
served as the regression test for the IPL-V interpreter. The ports of
Spice Lisp Flavors and LetS used in the KBEmacs runtime are by
g000001~\cite{g000001}. RAND Memorandum RM-3588-PR was obtained from
bitsavers. Listing excerpts from KBEmacs are copyright 1983/1984
Massachusetts Institute of Technology and are reproduced by permission
under the notice in those listings.

Finally, we acknowledge Herbert Simon, H.~S.~Kelly, Richard Waters,
Charles Rich, Kent Pitman, and the other members of the Programmer's
Apprentice project, whose work this paper merely re-runs.

{\footnotesize The transcriptions, the reanimation software, and
portions of the drafting of this paper were produced with the aid of
LLM-based tools, but all concepts and final editorial decisions were the
author's alone.}

\begin{thebibliography}{99}

\bibitem{backus1957fortran} J. W. Backus et al., ``The FORTRAN
  automatic coding system,'' in \textit{Proc. Western Joint
  Comput. Conf.}, 1957, pp. 188--198.

\bibitem{balzer1985} R. Balzer, ``A 15 year perspective on automatic
  programming,'' \textit{IEEE Trans. Softw. Eng.}, vol. SE-11, no. 11,
  pp. 1257--1268, Nov. 1985.

\bibitem{chen2021codex} M. Chen et al., ``Evaluating large language
  models trained on code,'' Jul. 2021, \textit{arXiv:2107.03374}.

\bibitem{g000001} g000001, ``spice-flavors'' and ``LetS,'' GitHub.
  %%% TODO verify access dates
  Accessed: Oct. 10, 2026. [Online]. Available:
  \url{https://github.com/g000001/spice-flavors},
  \url{https://github.com/g000001/LetS}

\bibitem{green1969} C. Green, ``Application of theorem proving to
  problem solving,'' in \textit{Proc. 1st Int. Joint Conf. Artif.
  Intell.}, 1969, pp. 219--239.
  %%% TODO verify pages

\bibitem{jimenez2024swebench} C. E. Jimenez, J. Yang, A. Wettig,
  S. Yao, K. Pei, O. Press, and K. Narasimhan, ``SWE-bench: Can
  language models resolve real-world GitHub issues?'' in
  \textit{Proc. Int. Conf. Learn. Represent. (ICLR)}, 2024.

\bibitem{lane2025eliza} R. Lane, A. Hay, A. Schwarz, D. M. Berry, and
  J. Shrager, ``ELIZA reanimated: Restoring the mother of all chatbots
  to one of the world's first time-sharing systems,'' \textit{IEEE
  Ann. Hist. Comput.}, vol. 47, no. 2, pp. 68--76, Apr.-Jun. 2025.

\bibitem{mannawaldinger1971} Z. Manna and R. J. Waldinger, ``Toward
  automatic program synthesis,'' \textit{Commun. ACM}, vol. 14, no. 3,
  pp. 151--165, Mar. 1971.

\bibitem{moewsshrager2026} D. Moews and J. Shrager, ``Executable
  archaeology: Reanimating the Logic Theorist,'' \textit{IEEE Ann. Hist.
  Comput.}, submitted, 2026.
  %%% TODO update status

\bibitem{newell1959gps} A. Newell, J. C. Shaw, and H. A. Simon,
  ``Report on a general problem-solving program,'' in
  \textit{Proc. Int. Conf. Inform. Process.}, Paris: UNESCO, 1960,
  pp. 256--264.

\bibitem{LT1957} A. Newell, J. C. Shaw, and H. A. Simon, ``Empirical
  explorations of the logic theory machine: A case study in
  heuristic,'' in \textit{Proc. Western Joint Comput. Conf.}, 1957,
  pp. 218--230.

\bibitem{newell1964ipl} A. Newell, F. M. Tonge, E. A. Feigenbaum,
  B. F. Green Jr., and G. H. Mealy,
  \textit{Information Processing Language-V Manual}. Englewood Cliffs,
  NJ, USA: Prentice-Hall, Inc., 2nd ed., 1964.

\bibitem{pitman1983} K. M. Pitman, ``Interfacing to the Programmer's
  Apprentice,'' MIT AI Lab., Cambridge, MA, USA, Working Paper 244,
  Feb. 1983.

\bibitem{pitmanscans} K. M. Pitman, ``netsettlement: kbe,'' GitHub
  (scans of \textit{Knowledge-Based Editor}, four-volume listing,
  Jan. 1985). Accessed: Oct. 13, 2026. [Online]. Available:
  \url{https://github.com/netsettler/netsettlement/tree/main/kbe}

\bibitem{richshrobe1976} C. Rich and H. E. Shrobe, ``Initial report on
  a LISP programmer's apprentice,'' MIT AI Lab., Cambridge, MA, USA,
  Tech. Rep. AI-TR-354, Dec. 1976. (Also \textit{IEEE Trans. Softw.
  Eng.}, vol. SE-4, no. 6, pp. 456--467, Nov. 1978.)

\bibitem{richwaters1987overview} C. Rich and R. C. Waters, ``The
  Programmer's Apprentice project: A research overview,'' MIT AI Lab.,
  Cambridge, MA, USA, AI Memo 1004, Nov. 1987.
  %%% TODO verify month

\bibitem{richwaters1988myths} C. Rich and R. C. Waters, ``Automatic
  programming: Myths and prospects,'' \textit{Computer}, vol. 21, no. 8,
  pp. 40--51, Aug. 1988.

\bibitem{richwaters1990book} C. Rich and R. C. Waters, \textit{The
  Programmer's Apprentice}. New York, NY, USA: ACM Press, and Reading,
  MA, USA: Addison-Wesley, 1990.

\bibitem{earepo} J. Shrager, ``Executable Archaeology,'' GitHub
  (HeuristicCoder, IPL-V, and LogicTheorist directories). Accessed:
  Oct. 16, 2026. [Online]. Available:
  \url{https://github.com/jeffshrager/ExecutableArchaeology}

\bibitem{kbemacsrepo} J. Shrager, ``kbemacs: MIT's KBEmacs from the
  January 1985 listings, running under SBCL,'' GitHub. Accessed:
  Oct. 16, 2026. [Online]. Available:
  \url{https://github.com/jeffshrager/kbemacs}

\bibitem{shrager2026} J. Shrager,
  ``Executable archaeology: Reanimating the Logic Theorist
  from its IPL-V source,''
  Mar. 13, 2026, \textit{arXiv:2603.13514v1}.

\bibitem{shragerrepo} J. Shrager. ``An IPL-V interpreter in Common Lisp
  that runs the original Logic Theorist,'' GitHub (archived; superseded
  by~\cite{earepo}). [Online]. Available:
  \url{https://github.com/jeffshrager/IPL-V}

\bibitem{simonP2349} H. A. Simon, ``Experiments with a heuristic
  compiler,'' RAND Corp., Santa Monica, CA, USA, Paper P-2349, Jun. 30,
  1961.

\bibitem{simonlisting1961} H. A. Simon, ``Heuristic Coder 7/16/61,''
  IPL-V program listing, Herbert A. Simon Papers, Box 14, Folder FF965,
  Carnegie Mellon Univ. Archives, Pittsburgh, PA, USA.

\bibitem{simon1963jacm} H. A. Simon, ``Experiments with a heuristic
  compiler,'' \textit{J. ACM}, vol. 10, no. 4, pp. 493--506, Oct. 1963.

\bibitem{simonRM3588} H. A. Simon, ``The heuristic compiler,'' RAND
  Corp., Santa Monica, CA, USA, Memorandum RM-3588-PR, May 1963.

\bibitem{waters1978} R. C. Waters, ``Automatic analysis of the logical
  structure of programs,'' Ph.D. dissertation, MIT, Cambridge, MA, USA,
  1978. (MIT AI Lab. Tech. Rep. AI-TR-492.)

\bibitem{waters1985tr} R. C. Waters, ``KBEmacs: A step toward the
  Programmer's Apprentice,'' MIT AI Lab., Cambridge, MA, USA, Tech. Rep.
  AI-TR-753, May 1985.

\bibitem{waters1985tse} R. C. Waters, ``The Programmer's Apprentice: A
  session with KBEmacs,'' \textit{IEEE Trans. Softw. Eng.}, vol. SE-11,
  no. 11, pp. 1296--1320, Nov. 1985.

\bibitem{waters1986aimag} R. C. Waters, ``KBEmacs: Where's the AI?''
  \textit{AI Mag.}, vol. 7, no. 1, pp. 47--56, Spring 1986.

\end{thebibliography}

\definecolor{darkblue}{RGB}{48, 107, 161}

\begin{IEEEbiography}{\color{darkblue} JEFF SHRAGER} is with Bennu Climate, Inc., Palo Alto, CA 94301 USA.  Contact them at jshrager@gmail.com.
\end{IEEEbiography}

\end{document}
